% fig_vl06_messung.tex — Was die Postulate 2, 4 und 5 zusammen festlegen:
% der Zustandsraum eines Qubits ist eine Scheibe, die Zeitentwicklung dreht
% darin, und der Detektor liest die waagerechte Koordinate ab.
%
% Chosen over the obvious "measurement cycle" flow chart because these
% postulates contain NO state-update rule -- there is no projection
% postulate in this lecture, so a "state after the measurement" panel would
% illustrate something the text does not claim.  What the text does claim
% has a shape, and the shape is nowhere drawn:
%
%   * Postulat 2 + Hermitezitaet + Spur eins give the lecture's
%         rho = ( p , x ; x^* , 1-p ),
%     and positivity is exactly det rho = p(1-p) - |x|^2 >= 0, i.e.
%         (p - 1/2)^2 + |x|^2 <= (1/2)^2 :
%     a disc of radius 1/2 about (1/2, 0).  Panel (a) is that inequality,
%     drawn in the slice Re x = 0 -- the plane the lecture's own
%     Hamiltonian preserves.  The boundary det rho = 0 is the pure states,
%     the centre is I/2, everything outside is a Hermitian trace-one
%     matrix that is NOT a state: hatched, not filled, so it survives
%     greyscale printing (and hatching is the right mark for "excluded").
%   * Postulat 5 with the lecture's H = ( 0,1 ; 1,0 ) and
%         U_t = cos(t/hbar) I - i sin(t/hbar) H
%     gives, for rho_t = U_t rho U_t^dag,
%         p_t  - 1/2 = r cos(2t/hbar + phi),   Im x_t = r sin(2t/hbar + phi),
%     with (p-1/2, Im x) = r(cos phi, sin phi).  So U_t is a rotation about
%     the centre of the disc at angular velocity 2/hbar: the radius r --
%     equivalently det rho -- is conserved.  The circles in (a) are those
%     orbits, not decoration.
%   * Postulat 4 with the lecture's detector F = |0><0| gives
%         tr(rho_t F) = p_t = 1/2 + r cos(2t/hbar + phi),
%     the horizontal coordinate of the rotating point.  Panel (b) is that
%     function.  Hence the point of the plate: the DISTANCE FROM THE CENTRE
%     in (a) is the AMPLITUDE of the measured frequency in (b).  The pure
%     state swings the full 0..1, the maximally mixed state answers 1/2
%     forever -- it is deaf to this experiment.
%
% The three marked states, all in the slice Re x = 0 (numbers used below):
%   |0><0|   r = 1/2, phi =  0     -> (p, Im x) = (1     , 0     )
%   rho      r = 0.3, phi = 70 deg -> (p, Im x) = (0.6026, 0.2819)
%   I/2      r = 0                 -> (p, Im x) = (0.5   , 0     )
% and their curves in (b) are 1/2 + r cos(2t/hbar + phi) with the same
% r and phi, so tr(rho F) at t = 0 is the p of the dot in (a) -- the
% lecture's own "Ansprechhaeufigkeit der Z-Messung" example.
%
% Colour encodes one thing only: the state, ordered by purity r
% (spacecadet r=0, munsell r=0.3, birrenorange r=1/2, the ramp of
% `tpclevels`).  A state has the same colour in both panels, and its orbit
% in (a) has the colour of the state that runs on it.  Construction lines
% (radius, projection onto p) are neutral grey.
%
% Both axes of (a) carry the same unit length (scale only axis with matched
% width/height), otherwise the circle -- the whole content of the panel --
% would render as an ellipse.

\begin{center}
\begin{tikzpicture}

% ─────────────────────────────────────────────────────────────
%  (a)  The state disc  det rho = p(1-p) - |x|^2 >= 0
% ─────────────────────────────────────────────────────────────
\begin{axis}[
    name=disc,
    tpcplot,
    scale only axis,
    width=4.61cm, height=5.11cm,     % 1.28 x 1.42 units at 3.6 cm/unit
    xmin=-0.14, xmax=1.14,
    ymin=-0.68, ymax=0.74,
    grid=none, clip=false, axis on top,
    xtick={0,0.5,1}, xticklabels={$0$,$\tfrac12$,$1$},
    ytick={-0.5,0,0.5}, yticklabels={$-\tfrac12$,$0$,$\tfrac12$},
    xlabel={$p$}, ylabel={$\imag x$},
    title={(a) \tpcifenglish{state disc}{Zustandsscheibe}},
]

  % Hermitesch, Spur eins, aber nicht positiv: kein Zustand.
  \fill[forbidden] (axis cs:-0.14,-0.68) rectangle (axis cs:1.14,0.74);

  % Das Innere der Scheibe freistellen: det rho >= 0.
  \addplot[fill=white, draw=none, domain=0:360, samples=181]
    ({0.5+0.5*cos(x)}, {0.5*sin(x)});

  % Bahn des gemischten Zustands unter U_t: det rho bleibt erhalten.
  \addplot[munsell, densely dotted, line width=0.8pt,
           domain=0:360, samples=181]
    ({0.5+0.3*cos(x)}, {0.3*sin(x)});
  \addplot[munsell, densely dotted, line width=0.8pt,
           -{Stealth[length=5pt]}, domain=118:152, samples=20]
    ({0.5+0.3*cos(x)}, {0.3*sin(x)});

  % Rand det rho = 0: die reinen Zustaende, ebenfalls eine Bahn.
  \addplot[birrenorange, line width=1.2pt, domain=0:360, samples=181]
    ({0.5+0.5*cos(x)}, {0.5*sin(x)});
  \addplot[birrenorange, line width=1.2pt, -{Stealth[length=5pt]},
           domain=175:205, samples=20]
    ({0.5+0.5*cos(x)}, {0.5*sin(x)});

  % Postulat 4: der Detektor liest die waagerechte Koordinate ab.
  \draw[-{Stealth[length=4pt]}, birrengray!75, line width=0.5pt, densely dotted]
    (axis cs:0.6026,0.2819) -- (axis cs:0.6026,-0.66);
  \node[lbl, font=\scriptsize, text=birrengray, anchor=north west,
        fill=white, fill opacity=0.9, text opacity=1, inner sep=1pt]
    at (axis cs:0.635,-0.52) {$\tr(\rho F)$};

  % Radius r = Abstand vom maximal gemischten Zustand.
  \draw[-{Stealth[length=5pt]}, birrengray, line width=0.7pt]
    (axis cs:0.5,0) -- (axis cs:0.6026,0.2819);
  \node[lbl, font=\scriptsize, text=birrengray, anchor=east]
    at (axis cs:0.525,0.175) {$r$};

  % Die drei Zustaende
  \node[circle, fill=birrenorange, inner sep=1.7pt] at (axis cs:1,0) {};
  \node[circle, fill=munsell,      inner sep=1.7pt] at (axis cs:0.6026,0.2819) {};
  \node[circle, fill=spacecadet,   inner sep=1.7pt] at (axis cs:0.5,0) {};

  \node[lbl, anchor=east]  at (axis cs:0.975,0) {$\ketbra{0}{0}$};
  \node[lbl, anchor=west]  at (axis cs:0.645,0.295) {$\rho$};
  \node[lbl, anchor=north east] at (axis cs:0.49,-0.05) {$\tfrac{\mathbb{I}}{2}$};

  % Rand und Umgebung benennen
  \draw[birrengray, line width=0.5pt] (axis cs:0.235,0.545) -- (axis cs:0.185,0.405);
  \node[lbl, font=\scriptsize, anchor=south west, inner sep=1pt]
    at (axis cs:-0.02,0.55) {\tpcifenglish{pure}{rein}: $\det\rho=0$};
  \node[lbl, font=\scriptsize, text=birrengray, anchor=south west,
        align=left, inner sep=1pt]
    at (axis cs:-0.10,-0.655) {$\det\rho<0$\\ \tpcifenglish{no state}{kein Zustand}};
  \node[lbl, font=\scriptsize, text=birrengray, anchor=north east]
    at (axis cs:1.14,0.72) {$\real x=0$};
  \node[lbl, font=\scriptsize, text=birrengray, anchor=north west]
    at (axis cs:0.125,0.315) {$U_t$};
\end{axis}

% ─────────────────────────────────────────────────────────────
%  (b)  Postulat 4 + 5:  tr(rho_t F) = 1/2 + r cos(2t/hbar + phi)
% ─────────────────────────────────────────────────────────────
\begin{axis}[
    tpcplot,
    at={(disc.south east)}, anchor=south west, xshift=15mm,
    scale only axis,
    width=5.7cm, height=5.11cm,
    xmin=0, xmax=6.2832, ymin=-0.06, ymax=1.10,
    clip=false,
    xtick={0,1.5708,3.1416,4.7124,6.2832},
    xticklabels={$0$,$\tfrac{\pi}{2}$,$\pi$,$\tfrac{3\pi}{2}$,$2\pi$},
    ytick={0,0.5,1}, yticklabels={$0$,$\tfrac12$,$1$},
    xlabel={$t/\hbar$},
    ylabel={$\tr(\rho_t F)$},
    every axis y label/.style={at={(axis description cs:-0.17,0.5)},
                               rotate=90, anchor=south, font=\small},
    title={(b) $Z$-\tpcifenglish{measurement}{Messung}, $F=\ketbra{0}{0}$},
]
  % tr(rho_t F) = 1/2 + r cos(2t/hbar + phi), Amplitude r = Radius aus (a)
  \addplot[birrenorange, line width=1.1pt, domain=0:6.2832, samples=241]
    {0.5+0.5*cos(deg(2*x))};
  \addplot[munsell, line width=1.1pt, domain=0:6.2832, samples=241]
    {0.5+0.3*cos(deg(2*x)+70)};
  \addplot[spacecadet, line width=1.1pt, domain=0:6.2832, samples=2]
    {0.5};

  \node[lbl, anchor=west] at (axis cs:6.40,1.00)  {$r=\tfrac12$};
  \node[lbl, anchor=west] at (axis cs:6.40,0.603) {$r=0.3$};
  \node[lbl, anchor=west] at (axis cs:6.40,0.50)  {$r=0$};
\end{axis}

\end{tikzpicture}
\end{center}
